% =====================================================================
% Section: Experiments -- Ablation study  (outline.md §3.5; defines sec:ablation,
%   referenced by 1_introduction.tex and 2_method.tex).
% Draft v1 (2026-06-17). Numbers are taken verbatim from the x4 ablation training
%   logs at iter 80k (paper/tables/table3_ablation_a1.tex / table4 / table5;
%   raw CSVs in paper/data/). Claims follow plan/review/method-experiment-
%   traceability.md ("消融实验的已知局限"): A1/A2 are reported as quality-neutral /
%   component-compatible (NOT incremental PSNR gains); A3 supports the balancing
%   MECHANISM (usage entropy), with the seed-variance sub-claim left as future work.
% Requires: booktabs (tables), pifont (Table IV \ding).
% =====================================================================
\subsection{Ablation study}
\label{sec:ablation}

We ablate the four design elements of DPR using the independent switches exposed in
Sec.~\ref{sec:method}. All ablations are conducted at scale $\times4$, following
CATANet~\cite{liu2025catanet} whose ablations are likewise performed at $\times4$, the
setting in which content routing on high-frequency texture is most exercised. To keep
the compute budget tractable, each variant is \emph{finetuned} from the converged
full-model checkpoint and trained for a short $80$k-iteration budget, with validation
on Set5 and Urban100 (Y channel, matched \texttt{crop\_border}) every $10$k iterations;
we report the iter-$80$k point. Besides PSNR/SSIM we log routing diagnostics per TAB
block---the normalized prototype-usage entropy $\mathcal{H}(u)$ of
Eq.~\eqref{eq:balance}, the active-prototype usage, and the learnable temperature---to
inspect the routing behavior directly.

\paragraph{Protocol caveat.}
Because every variant is finetuned from the \emph{same} full-model initialization for a
short budget, the disabled components are removed from an already converged
representation rather than learned from scratch. Under this protocol the reconstruction
differences among variants are small and fall within run-to-run training variance (the
full model's own validation PSNR fluctuates by $\sim0.03$\,dB between adjacent $10$k
checkpoints). We therefore read the PSNR/SSIM columns as evidence of
\emph{quality-neutrality and component compatibility}---each design element can be added
or removed without degrading reconstruction---rather than as cumulative accuracy gains,
and we substantiate the routing-level effects through the diagnostics. A fully separated
attribution of each switch would require retraining every variant from a
switch-neutral initialization to the CATANet $\times4$ from-scratch budget, which we
leave to future work.

\paragraph{A1: decoupled prototype confirmation.}
Table~\ref{tab:abl_a1} compares DPR with and without the query-refinement stage
(Eq.~\eqref{eq:proto_refine}). The two settings reach comparable reconstruction quality
(Set5 $32.55$ vs.\ $32.57$\,dB; Urban100 $26.84$ vs.\ $26.81$\,dB, both within the
variance band), confirming that adding the learnable query bank does not perturb
reconstruction. As the refinement is gated by a small-initialized scalar
(Eq.~\eqref{eq:proto_refine}), it acts as a residual confirmation step that keeps the
prototype slots stable while decoupling generation from confirmation (C2); we therefore
retain it by default. The motivation for dynamic, per-image prototypes over the
cross-batch EMA centers of the baseline (C1) is supported indirectly by the published
CATANet results rather than by this controlled switch, since re-implementing the EMA
center branch is outside the scope of our ablation.

\paragraph{A2: confidence-aware routing.}
Table~\ref{tab:abl_a2} adds the three confidence-aware components cumulatively:
confidence-aware ordering (Eq.~\eqref{eq:sortkey}), the IASA score gate
(Eq.~\eqref{eq:iasa_gate}), and the soft prototype fallback
(Eq.~\eqref{eq:soft_fallback}). Across the four configurations PSNR/SSIM stays within a
$\sim0.04$\,dB band on both Set5 and Urban100 and is not monotone in the number of
enabled components, which---given the shared-initialization protocol above---indicates
that the three switches are mutually compatible and quality-preserving rather than
individually accuracy-additive on top of a converged backbone. The components are kept
because they make the routing path interpretable and confidence-driven: the retained
per-token confidence $x^{\mathrm{s}}$ modulates ordering, the global gate, and the
fallback (C3), which we visualize through the $x^{\mathrm{s}}$ distributions in
Sec.~\ref{sec:routing_analysis}.

\paragraph{A3: prototype usage balancing.}
Table~\ref{tab:abl_a3} isolates the usage-balancing regularizer
(Eq.~\eqref{eq:balance}). Enabling it raises the block-averaged normalized usage entropy
from $0.6227$ to $0.6865$ on Urban100, and the increase is \emph{consistent across all
eight TAB blocks} ($b_0$--$b_7$; Table~\ref{tab:abl_a3_perblock}), i.e.\ the regularizer
drives a more balanced use of prototype slots and counteracts the collapse that the
learnable temperature would otherwise encourage (C4). This effect is visible even under
the short finetune protocol because it concerns the usage \emph{distribution} rather than
absolute reconstruction, and the PSNR/SSIM remains unchanged within variance. The
complementary part of C4---a reduction of seed-to-seed variance---requires repeating the
runs over multiple seeds; with a single seed available we report the balancing mechanism
here and leave the variance analysis as future work, and we do not claim a variance
reduction without those runs.

\paragraph{Summary.}
The ablations show that DPR's design elements are quality-neutral and mutually
compatible under a shared-initialization finetune, and that the balancing regularizer
produces a measurable, block-consistent increase in prototype-usage entropy. Stronger,
fully separated accuracy attributions are deferred to from-scratch retraining; the
present study supports the mechanism of each component without overstating its isolated
contribution to PSNR.
